%basic packages
\usepackage[utf8]{inputenc}
\usepackage[T1]{fontenc}
\usepackage{graphicx}
\usepackage[margin=1in]{geometry}
\usepackage[dvipsnames]{xcolor}
\usepackage{microtype}
\usepackage[skip=5pt, indent=40pt]{parskip}


%math
\usepackage{amsmath, amsthm, amsfonts, amssymb, mathtools, mathrsfs}
\usepackage{eucal}

%misc
\usepackage{float}
\usepackage[hyphens]{url}

\usepackage{hyperref}
\hypersetup{
    colorlinks,
    linkcolor={blue},
    urlcolor={blue},
    citecolor={red}
}
\usepackage{cleveref}

%%% macros

% general
\newcommand{\tildesim}{\mathbin{\sim}}
\renewcommand{\mod}[1]{\;(\text{mod}\;#1)}

% algebra
\DeclareMathOperator{\End}{End}
\DeclareMathOperator{\Aut}{Aut}
\DeclareMathOperator{\coker}{coker}
\DeclareMathOperator{\im}{im}
\DeclareMathOperator{\coim}{coim}
\renewcommand{\Im}{\operatorname{Im}}
\DeclareMathOperator{\Ker}{Ker}
\DeclareMathOperator{\Coker}{Coker}
\DeclareMathOperator{\Coim}{Coim}

%categories
\newcommand{\Set}{\mathcal{S}\mathrm{et}}
\newcommand{\Grp}{\mathcal{G}\mathrm{rp}}
\newcommand{\Top}{\mathcal{T}\mathrm{op}}
\newcommand{\Ring}{\mathcal{R}\mathrm{ing}}
\newcommand{\Cat}{\mathcal{C}\mathrm{at}}
\newcommand{\sSet}{\mathrm{s}\mathcal{S}\mathrm{et}}
\newcommand{\Ab}{\mathcal{A}\mathrm{b}}
\newcommand{\Vect}{\mathcal{V}\mathrm{ect}}
\newcommand{\Alg}{\mathcal{A}\mathrm{lg}}
\newcommand{\Ch}{\mathcal{C}\mathrm{h}}
\newcommand{\iCat}{\Cat_{\infty}}
\DeclareMathOperator{\Hom}{Hom}
\DeclareMathOperator{\Obj}{Obj}
\DeclareMathOperator{\Nat}{Nat}
\DeclareMathOperator{\Fun}{Fun}
\DeclareMathOperator{\Map}{Map}
\DeclareMathOperator{\Ho}{Ho}
\DeclareMathOperator{\Sing}{Sing}
\newcommand{\hh}{\mathrm{h}}
\newcommand{\Lan}{\mathrm{Ran}}
\newcommand{\Ran}{\mathrm{Ran}}
\DeclareMathOperator*{\hocolim}{hocolim}
\DeclareMathOperator*{\holim}{holim}


%theorems
\usepackage{thmtools}
\usepackage{tikz}
\usepackage{tikz-cd}
\usepackage{quiver}
\usepackage[framemethod=TikZ]{mdframed}
\mdfsetup{skipabove=1em,skipbelow=0em, innertopmargin=5pt, innerbottommargin=6pt}


\declaretheoremstyle[headfont=\bfseries, bodyfont=\normalfont]{definition}
\declaretheoremstyle[headfont=\bfseries, bodyfont=\normalfont, qed={$\blacksquare$}]{proofline}
\declaretheoremstyle[headfont=\bfseries, bodyfont=\normalfont, qed={\qedsymbol}]{egline}
\declaretheoremstyle[headfont=\bfseries, bodyfont=\normalfont]{dfn}
\declaretheoremstyle[headfont=\bfseries, bodyfont=\itshape]{thm}


\declaretheorem[numberwithin=subsection, style=definition, name=Definition]{definition}
\declaretheorem[sibling=definition, style=dfn, name=Remark]{remark}
\declaretheorem[sibling=definition, style=dfn, name=Example]{example}
\declaretheorem[sibling=definition, style=dfn, name=Notation]{notation}
\declaretheorem[sibling=definition, style=dfn, name=Exercise]{exercise}
\declaretheorem[sibling=definition, style=thm, name=Theorem]{theorem}
\declaretheorem[sibling=definition, style=thm, name=Corollary]{corollary}
\declaretheorem[sibling=definition, style=thm, name=Lemma]{lemma}
\declaretheorem[sibling=definition, style=dfn, name=Construction]{construction}
\declaretheorem[sibling=definition, style=thm, name=Proposition]{proposition}
\declaretheorem[sibling=definition, style=dfn, name=Terminology]{terminology}
\declaretheorem[sibling=definition, style=dfn, name=Intuition]{intuition}
\declaretheorem[sibling=definition, style=thm, name=Claim]{claim}

%solution environment
\declaretheorem[numbered=no, style=egline, name=Solution]{setsolution}
\newenvironment{solution}[1][]{\vspace{-10pt}\begin{setsolution}}{\end{setsolution}}
%proof environment
\declaretheorem[numbered=no, style=proofline, name=Proof]{replacementproof}
\renewenvironment{proof}[1][\proofname]{\begin{replacementproof}}{\end{replacementproof}}

\newtheorem*{note}{Note}
\newtheorem*{prev}{As Previously Seen}


%fancy headers
\usepackage{fancyhdr}
\pagestyle{fancy}
\fancyhead{}\fancyfoot{}
\fancyfoot[C]{\thepage}
\renewcommand{\headrulewidth}{0pt}

\author{Grant Talbert}
